\subsection{The coefficient residue when simplex zero has new zero-forms}

\paragraph{A-S09-17 (Proved; the residual map to $\mathbb Q/\mathbb Z$ and an actual two-cell obstruction).}
Fix one positive weight.  Let $X$ be a simplicial nonnegative
cochain complex with torsion-free terms and coefficient
$H^0(X)=\underline{\mathbb Z}$.  Let $f_i:X\to A_i$,
$i=0,1$, preserve coefficients, where both targets have
coefficient $H^0=\underline{\mathbb Z}$ and $H^1=0$.
Their terms may have torsion.  There is a canonical
simplicial homomorphism
\[
 \varepsilon:\operatorname{tors}H^1(X)
                         \longrightarrow\underline{\mathbb Q/\mathbb Z}.
 \tag{S09.66}
\]
Every extension of the prescribed pair into a complex
with coefficient $H^0$ and $H^1=0$ makes this map zero.
There is also an explicit simultaneous degree-zero
extension killing exactly its kernel on torsion
$H^1$.  Its degree-zero groups are free when $X^1$
is termwise free.  No preservation of horizontal
homology is asserted by this construction.

\begin{proof}
Put
\[
 N_n=\{a\in X_n^0\otimes\mathbb Q:da\in X_n^1\}.
 \tag{S09.67}
\]
For $a\in N_n$, choose an integral primitive $h_i$
of $f_i(da)$.  The rationalized difference
$f_{i,\mathbb Q}(a)-h_i$ is a unique rational
coefficient $r_i$.  Replacing $h_i$ changes $r_i$
by an integer, so
$\theta_i(a)=r_i\bmod\mathbb Z$ is well defined.
It is additive and simplicial: any discrepancy
between choices is precisely a closed integral
coefficient.  On $X_n^0$ it is zero, while on a
rational coefficient $r$ it is $r\bmod\mathbb Z$.
The difference $\theta_1-\theta_0$ therefore descends
through the actual isomorphism
\[
 N/(X^0+\underline{\mathbb Q})
       \xrightarrow{\ \cong\ }\operatorname{tors}H^1(X),
                 \qquad a\longmapsto[da].           \tag{S09.68}
\]
Surjectivity follows by writing $\ell c=db$ for a
torsion cycle and using $a=b/\ell$.  The kernel
is as stated because two rational primitives of
the same integral boundary differ by a rational
coefficient.  This proves (S09.66), with
$\varepsilon=\theta_1-\theta_0$.

Define $E^0=\ker\theta_0\cap\ker\theta_1\subset N$
and retain the positive terms of $X$.
For $a\in E^0$, the element $f_{i,\mathbb Q}(a)$
has a unique integral representative with differential
$f_i(da)$: existence follows by adjusting $h_i$ by
the integral coefficient $r_i$, and uniqueness
follows because a closed integral coefficient cannot
vanish upon rationalization.  These representatives
give actual additive simplicial extensions of the
two maps, even for targets with torsion.
The image of $dE^0$ modulo $dX^0$ is exactly
$\ker\varepsilon$.  Indeed subtracting a rational
coefficient from any $a\in N$ with
$\theta_0(a)=\theta_1(a)$ makes both residues zero.
The kernel of $d:E^0\to X^1$ is exactly the
integral coefficient group.  Thus $H^0(E)=\mathbb Z$
and the stated part of $H^1$ is killed.
If $X_n^1$ is free, then $dE_n^0$ is free, and
$0\to\mathbb Z\to E_n^0\to dE_n^0\to0$ splits,
proving freeness.

Finally, suppose $\ell c=db$ and the image of $c$
has a primitive $u$ in an extension $Y$ with
coefficient $H^0$.  Then $\ell u-b$ is the same
integral coefficient under both target maps.
Dividing this equality by $\ell$ shows
$\theta_0(b/\ell)=\theta_1(b/\ell)$.  This proves
necessity without requiring $X\to Y$ to be injective.
\end{proof}

Here is an explicit occurrence of a nonzero (S09.66)
after the two-cell construction AT14, even though
the old partial source has $H^1=0$ and is rationally
a coefficient resolution.  Put $p=3$ and use the
simplicial free module $L=\Phi/\underline{\mathbb Z}$
from (S09.53).  With $G$ the standard theory, set
\[
 X=G\otimes\Lambda(xL)\otimes\Gamma(yL),
 \quad |x|=1,\ |y|=2,\quad
 \operatorname{wt}(x)=\operatorname{wt}(y)=2,
 \quad dx=p y,
 \tag{S09.69}
\]
with the indicated linear differential for every
parameter in $L$.  There are actual strict PD maps
$f_0,f_1:X\to G$, equal to the identity on $G$, given by
\[
 f_0(x[a])=f_0(y[a])=f_1(y[a])=0,
 \qquad f_1(x[a_0,\ldots,a_n])
                    =T\sum_i a_i\,dt_i.
 \tag{S09.70}
\]
The last expression is well defined modulo the constant
vector, is simplicial, and is vertically closed.
Thus the prescribed differential is respected.

The source has free terms, $q\leq w$, $X_0=S$,
$H^0(X_n)=S$, and $H^1(X_n)=0$.  In every simplex,
its rationalization is a coefficient resolution:
choose a basis of $L_n$ and split it as tensor
products of the odd-to-even disks $dx=p y$,
whose positive-arity differentials have nonzero
integer coefficients and hence contract over
$\mathbb Q$.  The assertions about $H^0$ and $H^1$
also follow directly in degree one, where multiplication
by $p$ on the free $x$-module is injective, together
with the integral contraction of $G$.
All horizontal rows are acyclic for $w>2q$.
Indeed, an arity-$m$, form-$j$ term of the last two
factors has weight $2m$ and $m\leq j\leq2m$.
Every contributing term has $m\leq q$ and
$w-2m>q-j$ under this bound; its $G$ factor is
therefore an acyclic free simplicial module.
The integral tensor contraction proves the claim.

Let $M=D(1)$.  Its top normalized generator is in
degree one, and its bottom generator is its boundary.
Define simplicial maps by
\[
 \begin{aligned}
 z:M&\longrightarrow Z^2(X^{2,*}),&
 z_{\rm top}&=y[0,1],& z_{\rm bottom}&=0,\\
 b:M&\longrightarrow X^{2,1},&
 b_{\rm top}&=x[0,1],& b_{\rm bottom}&=0.
 \end{aligned}
 \tag{S09.71}
\]
They satisfy $db=pz$.  The AT14 extension adjoins
$u(M)$ in form one and $v(M)$ in form zero, both
of weight two, with
\[
                         du=z,\qquad dv=p u-b.
 \tag{S09.72}
\]
Denote it by $Y$.  In simplex zero it is exactly
\[
             Y_0=S\otimes\Gamma(v_0)\otimes\Lambda(u_0),
                        \qquad dv_0=p u_0.          \tag{S09.73}
\]
In particular $[u_0]\in H^1(Y_0)^2$ has exact order $p$.

The residue below is independent of every target
primitive choice in AT14.  For either extending
target map $g_i:Y\to G$, the map $g_i u:M\to G^1$
is closed because $f_i z=0$.  Projectivity of $M$
and vertical exactness give a simplicial primitive
$a_i:M\to G^0$.  It is unique because
$\operatorname{Hom}(D(1),\underline{S^2})=0$.
The maps $f_i b$ similarly have unique primitives
$k_i:M\to G^0$.  They are given exactly by
\[
 k_0=0,\qquad
 (k_1)_{\rm top}=Tt_1,\qquad
 (k_1)_{\rm bottom}=T^2=2T^{[2]}.                    \tag{S09.74}
\]
Here $d_1(Tt_1)=0$ and $d_0(Tt_1)=T^2$, which
checks the disk map, while its vertical differential
is $Tdt_1$.  The difference
$g_i v-p a_i+k_i$ is a simplicial coefficient map
from $D(1)$, so it is zero.  Therefore
\[
                          g_i v=p a_i-k_i.           \tag{S09.75}
\]
At the bottom vertex, each $a_i$ is an integral
coefficient, $g_i u_0=0$, and (S09.75) proves
\[
                  \varepsilon([u_0])=-\frac{2}{3}
                                      \quad\bmod\mathbb Z\ne0.
 \tag{S09.76}
\]
One may also take the completely explicit choices
$g_0u=g_1u=0$, $g_0v=0$, $g_1v=-k_1$;
the calculation proves that all other choices
have the same residue.

The obstruction persists in arbitrarily large weights.
Multiplying (S09.73) by $T^{[r]}$ gives the torsion
class $[T^{[r]}u_0]$ of weight $r+2$, and
\[
 \varepsilon([T^{[r]}u_0])
              =-\frac{(r+1)(r+2)}{3}\quad\bmod\mathbb Z.
 \tag{S09.77}
\]
The arity-one summand in (S09.73) verifies that this
class has order three; other arities do not contribute
boundaries to it.  For every $r$ divisible by three,
(S09.77) is nonzero.  Consequently no eventual
admissible extension of this $Y$ can preserve both
maps: at all sufficiently large such weights its
coefficient $H^0$ and vanishing $H^1$ would contradict
the necessity proved above.  This is an explicit
obstruction to the indicated repair of a partial
source.  Neither $X$ nor $Y$ is an admissible object,
and no counterexample to the main common-source
assertion is claimed.

The exact finite certificate
\texttt{code/A\_cosimplicial\_residue\_certificate.py}
checks the derivative and endpoint values in (S09.74),
and checks (S09.77) for $0\leq r\leq120$ using
integer arithmetic.  It also stores the integer
two-simplex differential and face matrices in
weights two through eight and verifies all $238$
columns of the integral Stokes identity.  The
output is
\texttt{logs/A\_cosimplicial\_residue\_certificate.json};
the measured runtime was $0.004414819995872676$ seconds.
The proof of the unbounded family is the formula
(S09.77) and its congruence for $r\equiv0\pmod3$,
not an extrapolation of those finite checks.
